\section{Using the repaired kernel to select a direction (NS-71)}
\label{sec:ns71-selected-preconditioner}

\paragraph{Question and scope.}
NS-70 excludes a uniform comparison over every coefficient vector.
It does not exclude using the repaired kernel to choose one direction,
then measuring that direction with the true complete Gram. The finite
prediction tested here is that such a choice captures more of the
optimized improvement than the unmodified correlation vector. A good
finite direction is a computational finding, not a cofinal estimate.

On the new block $N<n\leq2N$, let $D_N$ and $H_N$ denote the raw
and old-space-projected actual difference Grams. Let $D_{0,N}$ and
$H_{0,N}$ be the corresponding Grams for the repaired kernel $K_0$.
All four matrices are strictly positive definite. In defining $H_{0,N}$,
both the old projection and the new block use $K_0$ consistently;
this is not the actual Hilbert-space projection.
For the actual optimized correlation vector $g$, consider
\[
 v_{\rm raw}=D_{0,N}^{-1}g,\qquad
 v_{\rm projected}=H_{0,N}^{-1}g.
\]
Both are only rules for selecting coefficients. Every reported energy
uses $H_N$ or $D_N$, with the complete arithmetic remainder retained.

\begin{proposition}[Selected-direction improvement and the precise open input]
\label{prop:ns71-selected-input}
For either positive model matrix $A_N\in\{D_{0,N},H_{0,N}\}$ put
\[
 v=A_N^{-1}g,\quad m_N=g^TA_N^{-1}g,\quad
 C_N=\frac{v^TH_Nv}{v^TA_Nv},\quad
 C_N^{\rm raw}=\frac{v^TD_Nv}{v^TA_Nv}.
\]
Whenever $g\ne0$, all denominators are positive, and
\begin{equation}
 E_N-E_{2N}\geq Q_N(v)=\frac{m_N}{C_N}
                  \geq\frac{m_N}{C_N^{\rm raw}}>0.
 \label{eq:ns71-selected-ratio}
\end{equation}
If along $N_j=2^jN_0$ independently justified estimates give
\[
 \frac{m_{N_j}}{E_{N_j}}\geq\alpha_j>0,
 \qquad C_{N_j}\leq B_j,
 \qquad \sum_j\frac{\alpha_j}{B_j}=\infty,
\]
then $E_{N_j}\to0$ and RH follows by the fixed smoothed NB criterion.
The same statement holds using $C_{N_j}^{\rm raw}$ in place of $C_{N_j}$.
No such infinite sequence of estimates is proved here.
\end{proposition}
\begin{proof}
Since $A_Nv=g$, $v^Tg=v^TA_Nv=m_N$. Substitute in
\eqref{eq:ns69-selected-gain}. Orthogonal projection gives
$H_N\leq D_N$, proving the second inequality with all old-space
cross terms accounted for. The assumed bounds yield
$E_{2N_j}\leq(1-\alpha_j/B_j)E_{N_j}$.
The exact gain ensures $0<\alpha_j/B_j\leq1$; if equality occurs,
the error vanishes already. Otherwise the divergent sum makes the
product of contraction factors tend to zero. Monotonicity extends
convergence from the dyadic sequence to all $N$. The criterion in
Proposition~\ref{prop:ns61-rh-equivalence} then applies.
This uses a selected scalar ratio, not a uniform matrix comparison.
\end{proof}

For its geometric meaning, let $u=H_N^{-1/2}g$ and
$z=H_N^{1/2}v$. Then the exact fraction of the full improvement is
\[
 \frac{Q_N(v)}{E_N-E_{2N}}=
       \frac{(u^Tz)^2}{\|u\|_2^2\|z\|_2^2}.
\]
Thus the finite experiment measures an angle in the actual projected
energy, including the arithmetic terms. A Euclidean angle or a small
entrywise kernel remainder would not be a substitute.

\paragraph{Certified finite outcome.}
The table gives outward intervals for the fraction of the full
improvement $E_N-E_{2N}$, using the actual projected energy throughout.
\begin{center}\small
\begin{tabular}{rccc}
\toprule
$N$ & $v=g$ & $v=D_{0,N}^{-1}g$ & $v=H_{0,N}^{-1}g$\\
\midrule
16 & $(.2332,.2333)$ & $(.7167,.7169)$ & $(.8390,.8391)$\\
32 & $(.3258,.3259)$ & $(.8580,.8581)$ & $(.8252,.8254)$\\
64 & $(.2723,.2724)$ & $(.9593,.9594)$ & $(.9577,.9578)$\\
128 & $(.2946,.2948)$ & $(.9729,.9730)$ & $(.9781,.9782)$\\
256 & $(.3840,.3842)$ & $(.9597,.9598)$ & $(.9705,.9707)$\\
\bottomrule
\end{tabular}\end{center}
Both rules outperform $v=g$ at every checked size. The projected
rule captures more than $97\%$ of the full improvement at $N=128,256$;
neither efficiency is asserted monotone. At $N=256$ it reduces the
squared residual by a fraction in $(.08081,.08083)$, compared with
the full optimum in $(.08326,.08328)$. The corresponding numerator
and selected ratio in Proposition~\ref{prop:ns71-selected-input}
satisfy
\[
 .7668< m_{256}/E_{256}<.7669,\qquad
 9.4884<C_{256}<9.4886.
\]
These two finite bounds make explicit where the remaining arithmetic
problem sits; no size-independent bound is inferred.

All five sizes were evaluated at 256 and 384 bits with the full
Bernoulli/Hurwitz Gram tail. A doubled series cutoff at $N=256$
replays the largest problem. The verifier matches $115$ quantities
between precisions, $23$ at the doubled cutoff, and $16$ overlapping
quantities against the separate NS-69 implementation. The minimum
recorded accuracy is $88$ bits. All model and actual solves retain
their ball enclosures. The Maxima checks concern the exact quotient
and projection algebra, not the missing arithmetic estimates.

\paragraph{Stop condition.}
The model inversion is explicit and numerically useful on the checked
blocks. The arithmetic numerator $m_N/E_N$ and the selected true-energy
ratio still need cofinal estimates. Finite efficiency, by itself,
does not control either quantity as $N\to\infty$. NS-70's obstruction
to all-coefficient comparisons is respected throughout; no new
zero-free region, uniform Weil floor, G2 or RH proof follows.
